\documentclass[11pt]{article}

\usepackage[margin=1in]{geometry}
\usepackage{amsmath,amssymb,amsthm,mathtools}
\usepackage{microtype}
\usepackage{enumitem}
\usepackage[hidelinks]{hyperref}

\newtheorem{assumption}{Assumption}
\newtheorem{proposition}{Proposition}
\newtheorem{lemma}{Lemma}
\newtheorem{theorem}{Theorem}
\newtheorem{corollary}{Corollary}
\theoremstyle{remark}
\newtheorem{remark}{Remark}

\newcommand{\R}{\mathbb{R}}
\newcommand{\F}{\mathcal{F}}
\newcommand{\Nb}{\mathcal{N}}
\newcommand{\U}{\mathcal{U}}
\newcommand{\C}{\mathcal{C}}
\newcommand{\Ts}{\mathcal{T}_{\mathrm{s}}}
\newcommand{\Tg}{\mathcal{T}_{\mathrm{g}}}
\newcommand{\Steer}{\operatorname{Steer}_{\mathrm{CBF}}}
\newcommand{\Nearest}{\operatorname{Nearest}}
\newcommand{\interior}{\operatorname{int}}
\newcommand{\pospart}[1]{\left[#1\right]_{+}}
\newcommand{\rcert}{r_{\mathrm{cert}}}
\newcommand{\kreq}{\kappa_{\mathrm{req}}}
\newcommand{\unom}{u_{\mathrm{nom}}}
\newcommand{\kmax}{\kappa_{\max}}

\title{Probabilistic Completeness of CBF-Steered RRTConnect\\
\large A Local-Transparency Proof}
\author{}
\date{}

\begin{document}
\maketitle

\begin{abstract}
We prove probabilistic completeness for an RRTConnect planner whose local steering
primitive is filtered by control barrier function (CBF) constraints. The argument is
local: around any positive-clearance solution path there is a neighborhood of uniformly
positive radius in which sufficiently short nominal steering commands are CBF-feasible and
therefore traverse the quadratic-program filter unchanged. Inside that neighborhood CBF
steering coincides exactly with geometric steering, so the ball-covering argument for
geometric RRT applies verbatim. We further show that treating the CBF gain as a decision
variable---minimizing it lexicographically subject to a cap---does not change the control
the filter returns, so the adaptive-gain and fixed-gain planners are the same planner and
admit a single completeness theorem rather than a theorem and a corollary. Two hypotheses
are isolated as the ones an implementation can violate: the steering primitive must impose
no positive lower bound on extension length, and the safety certificate rather than the
constraint gradients must carry collision-freeness.
\end{abstract}

\section{Preliminaries}
\label{sec:prelim}

\subsection{Configuration space}

Let $Q\subset\R^d$ be a bounded configuration space. Unless stated otherwise $\|\cdot\|$
is the Euclidean norm and $\|\cdot\|_\infty$ the maximum norm; $B_r(q)$ is the open
Euclidean ball of radius $r$ about $q$, $\interior A$ the interior of $A$, $\mu$ Lebesgue
measure, and $[a]_+:=\max\{a,0\}$. A path is a continuous map $\pi:[0,1]\to Q$; its image
$\pi([0,1])$ is compact.

\subsection{Barrier-defined free space}

Let $h_i:Q\to\R$, $i=1,\dots,M$, be barrier functions encoding workspace collision,
self-collision, and any further validity constraints, and set
\begin{equation}
\F:=\{q\in\interior Q: h_i(q)>0,\ i=1,\dots,M\}.
\label{eq:free-space}
\end{equation}
The inequality is strict because the argument requires an interior solution path. A typical
workspace barrier is $h_i(q)=d(p_i(q))-r_i-m_i$ for a sphere center $p_i(q)$, a
signed-distance field $d$, a collision radius $r_i$ and a safety margin $m_i\ge 0$; self-collision
barriers are the analogous pairwise separations. Nothing below depends on these forms.
Where margins are used, $\F$ is understood throughout as the \emph{margin-inflated} free
space that the filter actually enforces, and Assumption~\ref{as:clearance} is a statement
about that set rather than about geometric contact.

We deliberately do \emph{not} assume $h_i\in C^1$. Interpolated grid distance fields, the
common implementation, are Lipschitz but only piecewise differentiable, and the proof does
not need more than the following.

\begin{assumption}[Filter data]
\label{as:filter-data}
For each $i$ the filter evaluates at $q$ a constraint vector $g_i(q)\in\R^d$ and the value
$h_i(q)$, and imposes the row $g_i(q)^\top u+\kappa h_i(q)\ge 0$. The map $q\mapsto h_i(q)$
is continuous and $g_i$ is bounded on bounded sets.
\end{assumption}

When $h_i$ is differentiable, $g_i=\nabla h_i$; the weaker statement is what the proof
consumes. Note that $g_i$ enters only through the question of whether a nominal command is
CBF-feasible. Collision-freeness is established separately, by the certificate of
Section~\ref{sec:certificate}, and never by the rows. This separation is what allows the
regularity hypothesis to be this weak.

\subsection{Steering, controls, and the CBF program}

$\Ts$ and $\Tg$ denote the start and goal trees, $\eta>0$ the maximum extension range, and
$x_{\mathrm{near}}=\Nearest(\mathcal T,x_{\mathrm{rand}})$ a nearest tree vertex. Let
$\Delta t>0$ be the integration step and $\bar v_k>0$ the componentwise speed limits. The
admissible control set is state dependent, since a step must respect the configuration
limits it ends in:
\begin{equation}
\U(q):=\bigl\{u\in\R^d:\ |u_k|\le\bar v_k\ \ \forall k,\ \ q+\Delta t\,u\in Q\bigr\}.
\label{eq:control-set}
\end{equation}
For a one-step extension toward $z$ the nominal command is $\unom=(z-q)/\Delta t$. We write
$\Steer(q,z)$ for the endpoint the CBF steering routine returns when asked to steer from
$q$ toward $z$.

Given a gain $\kappa\ge 0$ and a weight $W\succ 0$, the CBF quadratic program at $q$ is
\begin{equation}
\min_{u\in\U(q)}\ \tfrac12(u-\unom)^\top W(u-\unom)
\quad\text{s.t.}\quad
g_i(q)^\top u+\kappa h_i(q)\ge 0,\ \ i=1,\dots,M.
\label{eq:qp}
\end{equation}
For $q\in\F$ define the \emph{required gain} of a command $u$,
\begin{equation}
\kreq(q,u):=\max_{i=1,\dots,M}\pospart{-\frac{g_i(q)^\top u}{h_i(q)}},
\label{eq:kreq}
\end{equation}
which is well defined because $h_i(q)>0$ there.

\subsection{Probabilistic completeness}

Let $S_N$ be the event that the planner has returned a valid solution by the end of outer
iteration $N$. The planner is probabilistically complete on the class of instances
considered here if $\Pr(S_N)\to 1$ as $N\to\infty$ for every feasible instance satisfying
the hypotheses below. The only sampling property used is that every open ball in a
neighborhood of the reference path carries strictly positive sampling probability; a
sufficient condition is a sampling density bounded below by some $f_{\min}>0$ there.

\section{Setup}
\label{sec:setup}

\begin{assumption}[Positive-clearance solution]
\label{as:clearance}
There is a path $\pi:[0,1]\to\F$ with $\pi(0)=q_{\mathrm s}$ and $\pi(1)=q_{\mathrm g}$
whose image lies in the interior of every validity domain the collision representation
requires, including the domain of the distance field and the configuration limits.
\end{assumption}

\begin{lemma}[Compact tube]
\label{lem:tube}
Under Assumptions~\ref{as:filter-data} and~\ref{as:clearance} there exist $\delta_U>0$,
$\underline h>0$, $G<\infty$ and $b>0$ such that the open tube
\begin{equation}
\Nb:=\bigcup_{t\in[0,1]}B_{\delta_U}\bigl(\pi(t)\bigr)
\label{eq:tube}
\end{equation}
has compact closure $\overline\Nb\subset\F$ and satisfies, for all $q\in\overline\Nb$ and
all $i$,
\begin{equation}
h_i(q)\ge\underline h,
\qquad
\|g_i(q)\|\le G,
\qquad
\operatorname{dist}_\infty(q,\partial Q)\ge b.
\label{eq:tube-bounds}
\end{equation}
\end{lemma}

\begin{proof}
$\pi([0,1])$ is compact and each $h_i$ is continuous, so
$h_0:=\min\{h_i(\pi(t)):t\in[0,1],\,i\le M\}>0$. The set
$V:=\{q\in\interior Q: h_i(q)>h_0/2\ \forall i\}$ is open and contains $\pi([0,1])$, so by
compactness some $\delta_U>0$ makes the closed $\delta_U$-neighborhood of $\pi([0,1])$ a
subset of $V$. Take $\Nb$ as in \eqref{eq:tube} and $\underline h:=h_0/2$. Since $Q$ is
bounded, $\overline\Nb$ is compact; $G$ then exists by Assumption~\ref{as:filter-data} and
$b>0$ because $\overline\Nb$ is a compact subset of the open set $\interior Q$.
\end{proof}

Taking these bounds over the compact closure $\overline\Nb$ rather than over the open tube
is what makes them attained and hence uniform; this is used at every point below.

\section{The Certified Local Free Region}
\label{sec:certificate}

Suppose that for every $q\in\overline\Nb$ there are constants $\Lambda_{i,k}\ge0$, uniform
in $q$, with
\begin{equation}
\|p_i(q+\Delta q)-p_i(q)\|\ \le\ \sum_{k=1}^{d}\Lambda_{i,k}|\Delta q_k|=:A_i(\Delta q)
\qquad\text{for all }\Delta q,
\label{eq:lever}
\end{equation}
and let $s_i(q)>0$ be the clearance available to row $i$. For a distance field with
Lipschitz constant $L$ one may take $s_i(q)$ to be $h_i(q)/L$ intersected with the
remaining travel before $p_i(q)$ leaves the field's domain; for a self-collision row it is
the pairwise separation, with no Lipschitz factor. Define
\begin{equation}
\C(q):=\{\Delta q: A_i(\Delta q)\le s_i(q),\ i=1,\dots,M\},
\qquad
\rcert(q):=\min_i\frac{s_i(q)}{\sum_{k}\Lambda_{i,k}},
\label{eq:cert}
\end{equation}
rows with zero denominator being omitted from the minimum. Each $A_i$ is positively
homogeneous, $A_i(\tau\Delta q)=\tau A_i(\Delta q)$ for $\tau\in[0,1]$, so membership of
$\Delta q$ in $\C(q)$ certifies the whole segment $q+\tau\Delta q$ and not merely its
endpoint. Consequently
\begin{equation}
\|\Delta q\|_\infty<\rcert(q)
\quad\Longrightarrow\quad
[q,q+\Delta q]\subset\F.
\label{eq:cert-segment}
\end{equation}
Because the $\Lambda_{i,k}$ are uniform and $s_i$ is bounded below on $\overline\Nb$---by
$\underline h/L$ from Lemma~\ref{lem:tube} together with the domain clearance guaranteed by
Assumption~\ref{as:clearance}---the infimum
\begin{equation}
r_0:=\inf_{q\in\overline\Nb}\rcert(q)>0
\label{eq:r0}
\end{equation}
is strictly positive.

That \eqref{eq:cert-segment} bounds a whole segment, and does so without reference to any
$g_i$, is the reason the regularity hypothesis of Assumption~\ref{as:filter-data} suffices:
the rows decide feasibility, the certificate decides safety, and only the latter must hold
between sample points.

\section{The Gain Is Not a Degree of Freedom}
\label{sec:gain}

It is natural to treat $\kappa$ as a decision variable rather than a constant, preferring
among least-deviating controls the one demanding the smallest allowance to spend clearance.
Written out, with a cap $\kmax>0$, this is the lexicographic program
\begin{equation}
\operatorname*{lexmin}_{u,\kappa}\ \Bigl(\tfrac12(u-\unom)^\top W(u-\unom),\ \kappa\Bigr)
\quad\text{s.t.}\quad
g_i(q)^\top u+\kappa h_i(q)\ge 0\ \forall i,
\quad u\in\U(q),
\quad 0\le\kappa\le\kmax.
\label{eq:lex}
\end{equation}
The following proposition shows that \eqref{eq:lex} is not a different filter from
\eqref{eq:qp}, and it is what collapses the adaptive and fixed-gain cases into one theorem.

\begin{proposition}[Gain reduction]
\label{prop:reduction}
Let $q\in\F$, $W\succ0$ and $\kmax>0$, and suppose the feasible set of \eqref{eq:lex} is
nonempty. Then its optimal control $u^\ast$ is exactly the unique solution of the
fixed-gain program \eqref{eq:qp} at $\kappa=\kmax$, and the optimal gain is
$\kappa^\ast=\kreq(q,u^\ast)$.
\end{proposition}

\begin{proof}
Write $\mathrm{Feas}(\kappa):=\{u\in\U(q): g_i(q)^\top u+\kappa h_i(q)\ge0\ \forall i\}$.
Because $h_i(q)>0$ for $q\in\F$, row $i$ is equivalent to
$\kappa\ge-g_i(q)^\top u/h_i(q)$, so $\mathrm{Feas}$ is nondecreasing in $\kappa$ and a
control $u$ is feasible for \emph{some} $\kappa\in[0,\kmax]$ if and only if
$u\in\mathrm{Feas}(\kmax)$. The projection onto $u$ of the feasible set of \eqref{eq:lex}
is therefore $\mathrm{Feas}(\kmax)$, and the primary objective of \eqref{eq:lex} is
minimized over exactly the feasible set of \eqref{eq:qp} at $\kappa=\kmax$. That objective
is strictly convex because $W\succ0$, so its minimizer $u^\ast$ is unique. The secondary
objective then minimizes $\kappa$ over the singleton $\{u^\ast\}$, giving
$\kappa^\ast=\max\{0,\kreq(q,u^\ast)\}=\kreq(q,u^\ast)$.
\end{proof}

\begin{corollary}
\label{cor:same-planner}
The adaptive-gain filter \eqref{eq:lex} and the fixed-gain filter \eqref{eq:qp} with
$\kappa=\kmax$ return the same control at every configuration, hence induce the same
steering map $\Steer$ and the same planner. Any completeness statement for one is a
completeness statement for the other.
\end{corollary}

Accordingly we fix a gain $\kappa>0$ for the remainder, understanding it as $\kmax$ when
the lexicographic form is used. The minimized gain $\kappa^\ast$ retains an independent
purpose: it certifies that along the returned control every barrier obeys
$h_i(t)\ge h_i(0)e^{-\kappa^\ast t}$, an envelope tighter than the cap promises. It is a
certificate, not a control law.

\begin{lemma}[Transparency condition]
\label{lem:transparency-cond}
Let $q\in\F$ and suppose $\unom\in\U(q)$ and $\kreq(q,\unom)\le\kappa$. Then the solution
of \eqref{eq:qp} is $u^\ast=\unom$.
\end{lemma}

\begin{proof}
By \eqref{eq:kreq} and $h_i(q)>0$, the hypothesis $\kreq(q,\unom)\le\kappa$ states exactly
that $\unom$ satisfies every row at gain $\kappa$; with $\unom\in\U(q)$ it is feasible. The
objective then attains its global minimum value $0$, and since $W\succ0$ it does so only at
$u=\unom$.
\end{proof}

\begin{remark}[Computable surrogates]
\label{rem:surrogate}
Lemma~\ref{lem:transparency-cond} is monotone in its hypothesis: any
$\widehat\kappa\ge\kreq(q,\unom)$ with $\widehat\kappa\le\kappa$ certifies the same
conclusion. An implementation may therefore substitute any computable upper bound on
$\kreq$---for instance one read off the certificate region of
Section~\ref{sec:certificate}, which requires no $g_i$ at all---without affecting
soundness. A loose surrogate shrinks the radius $\rho_\kappa$ of Lemma~\ref{lem:transparency}
by the same factor and so degrades the rate constant of Theorem~\ref{thm:pc}, but cannot
destroy positivity and hence cannot affect completeness.
\end{remark}

\section{Local Transparency}
\label{sec:local}

\begin{lemma}[Uniform transparent neighborhood]
\label{lem:transparency}
There exists $\rho>0$ such that for every $q\in\Nb$ and every $z$ with $\|z-q\|<\rho$,
\begin{equation}
\Steer(q,z)=z,
\label{eq:transparent}
\end{equation}
and the entire steering segment $[q,z]$ lies in $\F$.
\end{lemma}

\begin{proof}
Put $\delta:=z-q$ and $\bar v_{\min}:=\min_k\bar v_k>0$, and set
\begin{equation}
\rho:=\tfrac12\min\Bigl\{\,r_0,\ \ \Delta t\,\bar v_{\min},\ \ b,\ \
\frac{\kappa\,\Delta t\,\underline h}{G}\,\Bigr\}>0,
\label{eq:rho}
\end{equation}
which is positive because each of the four quantities is, by \eqref{eq:r0} and
Lemma~\ref{lem:tube}. Assume $\|\delta\|<\rho$; we verify the three hypotheses of
Lemma~\ref{lem:transparency-cond} in turn and then the segment claim.

\emph{One step, within the speed box.} Since $\|\delta\|_\infty\le\|\delta\|<\Delta t\,\bar
v_{\min}\le\Delta t\,\bar v_k$ for every $k$, the request is met in a single integration
step and the nominal command $\unom=\delta/\Delta t$ satisfies $|(\unom)_k|<\bar v_k$.

\emph{Admissibility.} $\|\delta\|_\infty<b$ and $q\in\Nb$ give
$z=q+\Delta t\,\unom\in\interior Q$ by \eqref{eq:tube-bounds}, so $\unom\in\U(q)$.

\emph{Gain headroom.} By Cauchy--Schwarz and the bounds of Lemma~\ref{lem:tube},
\begin{equation}
\kreq(q,\unom)
=\max_i\pospart{-\frac{g_i(q)^\top\delta}{\Delta t\,h_i(q)}}
\le\frac{G}{\Delta t\,\underline h}\,\|\delta\|
<\frac{G}{\Delta t\,\underline h}\cdot\frac{\kappa\,\Delta t\,\underline h}{G}
=\kappa .
\label{eq:kreq-bound}
\end{equation}
Lemma~\ref{lem:transparency-cond} therefore applies and yields $u^\ast=\unom$, so the
executed step is
\[
q^+=q+\Delta t\,u^\ast=q+\Delta t\,\frac{z-q}{\Delta t}=z .
\]

\emph{Safety of the segment.} Finally $\|\delta\|_\infty\le\|\delta\|<r_0\le\rcert(q)$, so
\eqref{eq:cert-segment} gives $[q,z]\subset\F$.
\end{proof}

\begin{remark}[No lower bound on extension length]
\label{rem:no-floor}
Lemma~\ref{lem:transparency} asserts transparency on a \emph{ball}, not an annulus: the
conclusion must hold for arbitrarily short requests. This is a genuine obligation on the
steering primitive and not a formality, because
Theorem~\ref{thm:pc} covers the path in balls and requires the tree to be extendable to a
sample lying anywhere inside one, however close that sample falls to the vertex it grows
from. Implementations that abandon a rollout whose applied control is small in magnitude
violate it, since a short request has a small nominal command by construction---$\unom$
scales with $\|z-q\|$---so such a test refuses precisely the short extensions the argument
depends on, and imposes a positive floor on extension length. Any work limit of this kind
must therefore be conditioned on the step failing to complete the request, not on the
magnitude of the command alone.
\end{remark}

\begin{remark}
Lemma~\ref{lem:transparency} does not claim the CBF steer is globally straight. Long
extensions, or extensions near $\partial\F$, may be deflected, truncated, or refused.
Completeness needs only a uniformly positive neighborhood of a positive-clearance path in
which exact steering is recovered.
\end{remark}

\section{Probabilistic Completeness}
\label{sec:pc}

\begin{theorem}[Probabilistic completeness of CBF-RRTConnect]
\label{thm:pc}
Let Assumptions~\ref{as:filter-data} and~\ref{as:clearance} hold, together with the
certificate hypothesis \eqref{eq:lever}, and suppose
\begin{enumerate}[label=(\roman*)]
\item samples are drawn independently from a distribution whose density is bounded below by
      $f_{\min}>0$ on $\Nb$;
\item RRTConnect has maximum extension range $\eta>0$;
\item the goal tree contains $q_{\mathrm g}$;
\item every local extension of either tree uses $\Steer$, with a fixed gain $\kappa>0$ or,
      equivalently by Corollary~\ref{cor:same-planner}, with the lexicographic gain capped
      at $\kmax=\kappa$.
\end{enumerate}
Then $\Pr(S_N)\to1$ as $N\to\infty$.
\end{theorem}

\begin{proof}
Let $\rho>0$ be the radius of Lemma~\ref{lem:transparency} and choose
\begin{equation}
0<\nu<\min\Bigl\{\eta,\ \delta_U,\ \tfrac53\rho\Bigr\}.
\label{eq:nu}
\end{equation}

\emph{Waypoints.} $\pi$ is continuous on the compact interval $[0,1]$, hence uniformly
continuous, so there is $\omega>0$ with $\|\pi(t)-\pi(t')\|\le\nu/5$ whenever
$|t-t'|\le\omega$. Put $m:=\lceil1/\omega\rceil$ and $x_j:=\pi(j/m)$ for $j=0,\dots,m$.
Then $x_0=q_{\mathrm s}$, $x_m=q_{\mathrm g}$ and $\|x_{j+1}-x_j\|\le\nu/5$. (Uniform
continuity is used here in place of rectifiability, which a merely continuous path need not
possess.) Set $B_j:=B_{\nu/5}(x_j)$.

\emph{Advancement.} Suppose $\Ts$ contains some $x_j'\in B_j$ and the next sample satisfies
$x_{\mathrm{rand}}\in B_{j+1}$. Let $x_{\mathrm{near}}$ be the tree vertex nearest
$x_{\mathrm{rand}}$. Nearest-neighbor optimality and the triangle inequality give
\begin{equation}
\|x_{\mathrm{near}}-x_{\mathrm{rand}}\|
\le\|x_j'-x_{\mathrm{rand}}\|
\le\|x_j'-x_j\|+\|x_j-x_{j+1}\|+\|x_{j+1}-x_{\mathrm{rand}}\|
<\tfrac{3\nu}{5}<\rho,
\label{eq:three-fifths}
\end{equation}
where the last inequality is \eqref{eq:nu}. Since $3\nu/5<\nu<\eta$ the extension is not
truncated by the range. Moreover
\begin{equation}
\|x_{\mathrm{near}}-x_j\|
\le\|x_{\mathrm{near}}-x_{\mathrm{rand}}\|+\|x_{\mathrm{rand}}-x_j\|
<\tfrac{3\nu}{5}+\tfrac{2\nu}{5}=\nu<\delta_U,
\label{eq:near-in-tube}
\end{equation}
so $x_{\mathrm{near}}\in\Nb$. Lemma~\ref{lem:transparency} thus applies to the pair
$(x_{\mathrm{near}},x_{\mathrm{rand}})$ and yields
$\Steer(x_{\mathrm{near}},x_{\mathrm{rand}})=x_{\mathrm{rand}}$ along a segment in $\F$.
The sampled state is therefore inserted, and
\begin{equation}
\Ts\cap B_j\neq\varnothing\ \wedge\ x_{\mathrm{rand}}\in B_{j+1}
\quad\Longrightarrow\quad
\Ts^{\mathrm{new}}\cap B_{j+1}\neq\varnothing,
\label{eq:advance}
\end{equation}
which is the local advancement property of the geometric RRT proof. We note that
\eqref{eq:three-fifths} uses only that the nearest vertex is no farther than $x_j'$, so the
property is unaffected if the filter also inserts additional vertices elsewhere---for
instance the endpoints of deflected or truncated extensions---since a nearest-neighbor
distance over a superset of vertices can only decrease.

\emph{Iteration count.} Each $B_j\subset\Nb$ is an open ball of positive volume, so
$\Pr[x_{\mathrm{rand}}\in B_j]\ge f_{\min}\mu(B_{\nu/5})=:p>0$, and only the finite number
$m$ of advancements is required. Tree vertices are never removed, so once $\Ts$ meets $B_j$
it continues to. RRTConnect alternates which tree receives the random extension, so after
$N$ outer iterations $\Ts$ has had at least $n_N:=\lfloor N/2\rfloor$ opportunities, and by
\eqref{eq:advance} each succeeds with conditional probability at least $p$. Stochastic
domination by a binomial then gives
\begin{equation}
\Pr\bigl[\Ts\text{ has not reached }B_m\bigr]
\le\sum_{j=0}^{m-1}\binom{n_N}{j}p^{\,j}(1-p)^{\,n_N-j}\xrightarrow[N\to\infty]{}0
\label{eq:binomial}
\end{equation}
for fixed $m$ and $p>0$.

\emph{Connection.} The goal tree contains $q_{\mathrm g}=x_m$. Once $\Ts$ acquires a vertex
$y\in B_m$ we have $\|y-q_{\mathrm g}\|<\nu/5<\rho$, and $y\in\Nb$ by
\eqref{eq:near-in-tube}. Lemma~\ref{lem:transparency} is stated for every $q\in\Nb$ and
every target within $\rho$, so it applies whichever of the two endpoints the
implementation steers \emph{from}; this matters because RRTConnect validates goal-tree
edges in the reverse direction, rolling out from the sampled state toward the tree vertex.
Either orientation yields exact arrival, the trees connect, and a solution is returned.
Combining with \eqref{eq:binomial} gives $\Pr(S_N)\to1$.
\end{proof}

\begin{corollary}[Adaptive gain]
\label{cor:adaptive}
Under the hypotheses of Theorem~\ref{thm:pc}, the planner steered by the lexicographic
program \eqref{eq:lex} with cap $\kmax>0$ is probabilistically complete.
\end{corollary}

\begin{proof}
Immediate from Corollary~\ref{cor:same-planner} and Theorem~\ref{thm:pc} with
$\kappa=\kmax$.
\end{proof}

\begin{remark}[Role of the gain]
\label{rem:gain-role}
Decreasing $\kappa$ shrinks $\rho_\kappa=\kappa\Delta t\,\underline h/G$ and hence, through
\eqref{eq:rho} and \eqref{eq:nu}, the favorable sampling probability $p$; it never
eliminates the radius. The gain therefore governs the convergence rate this sufficient
condition certifies, not completeness itself. By Corollary~\ref{cor:same-planner} an
adaptive gain does not enlarge the set of preserved nominal commands relative to a fixed
gain at the same cap, so it is neither required for the guarantee nor able to improve it;
its value lies in the tighter decay envelope it certifies after the fact.
\end{remark}

\begin{remark}[Rate constants]
\label{rem:rate}
The bound \eqref{eq:binomial} is honest but severe. In a typical manipulator instance the
binding term of \eqref{eq:rho} is the one-step speed bound $\Delta t\,\bar v_{\min}$, some
two orders below the planner range $\eta$, and $p$ scales as $(\rho/\operatorname{diam}Q)^d$.
The theorem should accordingly be read as establishing that the guarantee is not forfeited
by the filter, with empirical evidence carrying the practical rate.
\end{remark}

\section*{Reference}

Kleinbort, M., Solovey, K., Littlefield, Z., Bekris, K.~E., and Halperin, D.
``Probabilistic Completeness of RRT for Geometric and Kinodynamic Planning with Forward
Propagation.'' arXiv:1809.07051.

\end{document}
